% pdftexja standalone: vertical document (shipout rotation)
\documentclass{article}
\usepackage[tate]{pdftexja}
% a vertical page geometry: \textwidth is the line length (vertical)
\setlength{\textwidth}{150mm}\setlength{\textheight}{100mm}
\setlength{\oddsidemargin}{0pt}\setlength{\topmargin}{0pt}
\setlength{\headheight}{0pt}\setlength{\headsep}{0pt}
\begin{document}
\ptjiftdir{\typeout{PTJTEST-PASS: tdir}}{\typeout{PTJTEST-FAIL: tdir}}
\ptjiftatepage{\typeout{PTJTEST-PASS: tatepage}}{\typeout{PTJTEST-FAIL: tatepage}}
\edef\x{\the\dimexpr\ptjcharheight+\ptjchardepth\relax}
\ifdim\x=10pt \typeout{PTJTEST-PASS: char cell 1zw}\else\typeout{PTJTEST-FAIL: char cell \x}\fi
\setbox0\hbox{\ptjtatechuyoko{12}}
\ifdim\wd0=10pt \typeout{PTJTEST-PASS: tcy advance}\else\typeout{PTJTEST-FAIL: tcy advance \the\wd0}\fi
これは縦組みの文章です。「括弧」と、句読点。数字\ptjtatechuyoko{12}も入る。
\ptjyokopage
\ptjiftatepage{\typeout{PTJTEST-FAIL: yokopage}}{\typeout{PTJTEST-PASS: yokopage}}
横組みのページ。
\ptjtatepage
縦組みに戻る。
\end{document}
